\documentclass[11pt]{article}
\usepackage[no-math]{fontspec}
\usepackage[french]{babel}
\usepackage{amsmath}
\usepackage{amssymb}
\usepackage{amsthm}
\usepackage{mathrsfs}
\usepackage{tikz-cd}
\usepackage{enumitem}
\usetikzlibrary{babel}

\newcommand{\printedpage}[1]{\marginpar{\footnotesize p.~#1}}
\newcommand{\msp}[1]{\marginpar{\footnotesize [#1]}}

\begin{document}

\printedpage{24}
% head: A. Grothendieck : Dérivateurs, Ch. VI, Édition des Universités de Montpellier II et Paris VII

\msp{page 41}

\begin{flushright}
15.11.90
\end{flushright}

Je crois qu j'ai enfin compris, cette nuit, la situation exacte. Il y a lieu, dans Cat$/S$ et dans les autres catégories du diagramme p. 8 (s'envoyant toutes dans Cat$/S$) d'envisager quatre sortes de ‘$W$-équivalences’. % [sic] qu j'ai -> que j'ai, letter missing

\begin{itemize}
\item[a)] La \textit{$W$-équivalence ‘grossière’} : \begin{tikzcd}[baseline=(s.base), row sep=small, column sep=small]
f : X \arrow[rr] \arrow[dr] & & Y \arrow[dl] \\
& |[alias=s]| S &
\end{tikzcd} dans Cat$/S$ est une $W$-équivalence grossière si en tant que flèche de Cat, c'est une $W$-équivalence. (\footnote[17]{stable par changement de base à la fois lisse et propre, plus généralement si $S' \longrightarrow S$ est une Hot-fibration.}).

\item[b)] La \textit{$W$-équivalence ‘par fibres’} (sur $S$) : $\forall\ s \in S$, $f_s : X_s \longrightarrow Y_s$ est une $W$-équivalence. (\footnote[18]{stable par tout chagement de base $S' \longrightarrow S$}). % [sic] chagement -> changement, letter missing

\item[c)] La \textit{$W$-équivalence ‘locale sur $S$’} : $\forall\ s \in S$, $f/s : X/s \longrightarrow Y/s$ est une $W$-équivalence. Signifie que c'est une $W$-équivalence, et le reste par tout changement de base $S' \longrightarrow S$ qui est un isomorphisme local. (\footnote[19]{devrait être stable par changement de base \textit{lisse} $S' \longrightarrow S$}).

\item[c$'$)] La \textit{$W$-équivalence ‘colocale sur $S$’} : $\forall\ s \in S$, $s\backslash f : s\backslash X \longrightarrow s\backslash Y$ est une $W$-équivalence. Signifie que c'est une $W$-équivalence, et le reste par tout changement de base $S' \longrightarrow S$ qui est un co-isomorphisme local. (\footnote[20]{devrait être stable par changement de base \textit{propre} sur $S$}).
\end{itemize}

Je vais résumer dans un énoncé récapitulatif les relations entre ces quatre notions de $W$-équivalence, en pensant au cas $W = W_\infty$. Une fois la chose faite, je me préoccuperai de dégager les conditions minimales sur un $W$ général qui assurent que c'est vrai pour $W$.

\textbf{Proposition 1.} \textit{Soit} \begin{tikzcd}[baseline=(s.base), row sep=small, column sep=small]
f : X \arrow[rr] \arrow[dr] & & Y \arrow[dl] \\
& |[alias=s]| S &
\end{tikzcd} \textit{dans} Cat$/S$.

\begin{itemize}
\item[1)] \textit{La $W$-équivalence locale sur $S$, et la $W$-équivalence colocale sur $S$, impliquent la $W$-équivalence grossière :}
\[
\begin{tikzcd}
\text{c} \arrow[dr, Rightarrow] & & \text{c}' \arrow[dl, Rightarrow] \\
& \text{a} &
\end{tikzcd}
\]

\printedpage{25}
% head: A. Grothendieck : Dérivateurs, Ch. VI, Édition des Universités de Montpellier II et Paris VII

\item[2)] \textit{Si $X$, $Y$ sont propres sur $S$, ou $X$, $Y$ lisses sur $S$, ou $X$, $Y$ des} Hot-\textit{fibrés sur $S$, alors la $W$-équivalence par fibres implique que c'est une $W$-équivalence} universelle \textit{sur $S$ ($W$-équivalence, et le reste par} tout \textit{changement de base $S' \longrightarrow S$), donc dans ces trois cas, on a}
\[
\begin{tikzcd}
& \text{b} \arrow[dl, Rightarrow] \arrow[dr, Rightarrow] & \\
\text{c} \arrow[dr, Rightarrow] & & \text{c}' \arrow[dl, Rightarrow] \\
& \text{a} &
\end{tikzcd}
\]
(\footnote[21]{Il suffit $X$, $Y$ $W$-propres sur $S$, ou $W$-lisses sur $S$, ou $W$-fibrés sur $S$.}).

Mais on peut préciser chacun de ces trois cas :

\item[(2 fib)] \textit{Si $X$, $Y$ sont} Hot-\textit{fibrés sur $S$, alors la $W$-équivalence grossière de $f$ implique que c'est une $W$-équivalence par fibres, donc les quatre notions de $W$-équivalence sont équivalentes pour $f$.}
\[
\text{a} \Longleftrightarrow \text{c} \Longleftrightarrow \text{c}' \Longleftrightarrow \text{b}
\]
(\footnote[22]{a $\Longrightarrow$ b seulement moyennant L5, voir ci-dessous p. 43.})

\item[(2 lisse)] \msp{page 42}\textit{Si $X$, $Y$ sont lisse sur $S$, alors la $W$-équivalence par fibres sur $S$ équivaut à la $W$-équivalence colocale sur $S$ et implique la $W$-équivalence} universelle \textit{sur $S$, donc dans ce cas} % [sic] lisse -> lisses, plural agreement with X, Y
\[
\begin{tikzcd}
\text{b} \arrow[r, Leftrightarrow] \arrow[d, Rightarrow] & \text{c}' \\
\text{c} \arrow[d, Rightarrow] & \\
\text{a} &
\end{tikzcd}
\]

\printedpage{26}
% head: A. Grothendieck : Dérivateurs, Ch. VI, Édition des Universités de Montpellier II et Paris VII
(\footnote[23]{il suffit même $X$, $Y$ $W$-lisses sur $S$.}).

\item[(2 prop)] \textit{Si $X$, $Y$ sont propres sur $S$, alors la $W$-équivalence par fibres sur $S$ équivaut à la $W$-équivalence locale sur $S$ :}
\[
\begin{tikzcd}
\text{b} \arrow[r, Leftrightarrow] \arrow[d, Rightarrow] & \text{c} \\
\text{c}' \arrow[d, Rightarrow] & \\
\text{a} &
\end{tikzcd}
\]
(\footnote[24]{Il suffit même $X$, $Y$ $W$-propres sur $S$.}).
\end{itemize}

\textsc{Démonstration} 1) n'est autre que L4 (\footnote[25]{Non, il y a confusion, L4 est plus faible, cf. p. 44 l'axiome L4 bis.}) et son dual L4$'$ pour un localiseur fondamental $W$, et on a vu que si on suppose L1, L3, alors L4 équivaut à L4$'$, et dans la notion de localiseur fondamental on supposera toujours que c'est vérifié. Également, nous supposons que L3 est vrai même sous la forme plus forte L3 quater (p. 1) : $W_\infty \subseteq W$, i.e. toute Hot-équivalence est une $W$-équivalence.

(2 lisse) Si $X$, $Y$ \texttt{[sont]} lisses sur $S$, alors les $X_s \longrightarrow s\backslash X$, $Y_s \longrightarrow s\backslash Y$ sont des $W_\infty$-équivalences, donc des $W$-équivalences (il suffirait même que $X$, $Y$ soient $W$-lisses$/S$), donc le diagramme commutatif
\[
\begin{tikzcd}
X_s \arrow[r, "f_s"] \arrow[d, "i_s^X"'] & Y_s \arrow[d, "i_s^Y"] \\
s\backslash X \arrow[r, "s\backslash f"] & s\backslash Y
\end{tikzcd}
\]
nous montre que $f_s$ est une $W$-équivalence si et seulement si $s\backslash f$ l'est, d'où b $\Longleftrightarrow$ c$'$. Pour voir que la $W$-équivalence par fibres est même une $W$-équivalence \textit{universelle} sur $S$, on est ramené à montrer que c'est une $W$-équivalence grossière (puisque la lissité \texttt{[est]} stable par changement de base sur $S$), et on gagne puisque c$'$ $\Longrightarrow$ a par 1).

Le cas (2 propre) est dual.

(2 fib) Si $X$, $Y$ sont des Hot-fibrés sur $S$, alors
\printedpage{27}
% head: A. Grothendieck : Dérivateurs, Ch. VI, Édition des Universités de Montpellier II et Paris VII
\msp{page 43}
les inclusions
\[
X_s \longrightarrow s\backslash X, \qquad X_s \longrightarrow X/s,
\]
et de même pour $Y$, sont des $W_\infty$-équivalences, donc des $W$-équivalences (il suffirait même que $X$, $Y$ soient des $W$-fibrations sur $S$ (\footnote[26]{comparer déf. 3, p. 117, dans (I).}), i.e. par définition que les foncteurs précédents soient des $W$-équivalences, et que la même condition soit satisfaite après tout changement de base sur $S$), l'argument dans le cas lisse resp. propre montre déjà que
\[
\begin{tikzcd}
\text{c} & \text{b} \arrow[l, Rightarrow] \arrow[r, Leftrightarrow] \arrow[d, Rightarrow] & \text{c}' \\
& \text{a,} &
\end{tikzcd}
\]
et que la $W$-équivalence par fibres (ou la $W$-équivalence locale sur $S$, resp. colocale sur $S$) implique déjà la $W$-équivalence \textit{universelle} sur $S$. (\footnote[27]{valable en supposant seulement $X$, $Y$ des $W$-fibrés sur $S$.})

La question a $\Longrightarrow$ b est par contre nettement plus délicate, et on a vu que ce n'est nullement le cas pour tout localisateur fondamental - sans doute pas pour celui associé au dérivateur Der($\mathbf{Z}$-Mod), en tous cas pas à celui associé à Der($k$-Mod), quand 2 est inversible dans $k$. (Il suffirait même que $\exists\ n \in \mathbf{N}$, $n \geq 2$, tel que $n$ inversible dans $k$ - le contrexemple typique est $X = S^1 \xrightarrow{\ n\ } S^1 = Y = S$, le revêtement étale d'ordre $n$ du cercle.)

L'implication a $\Longrightarrow$ b est donc une condition très forte sur $W$, je vais la noter L5 (l'ancienne L5, cf. page 2, ne méritant pas de survivre, car c'est une côte mal taillée).

\textbf{L5} \textit{Soit} \begin{tikzcd}[baseline=(s.base), row sep=small, column sep=small]
X \arrow[rr, "f"] \arrow[dr] & & Y \arrow[dl] \\
& |[alias=s]| S &
\end{tikzcd} \textit{dans} Cat\textit{, avec $X$, $Y$ des} Hot-\textit{fibrés sur $S$. Si $f$ est une $W$-équivalence (‘grossière’ par rapport à $S$), alors c'est une $W$-équivalence par fibres} (\footnote[28]{Dans XII (p. 12) on voit que si $W$ satisfait à L1, L3 quater, L4 bis (cf. page suivante), L5, alors \fbox{$W = W_\infty$}. Donc c'est parfaitement illusoire, si on suppose L5 de continuer à traîner un $W$.}).

On a vu qu'à un $\varepsilon$ près (‘lemme’ p. 23), L5 est une conséquence de

\textbf{L6} \textit{Soit $f : X \longrightarrow Y$ une $W$-équivalence, $Y' \longrightarrow Y$ un morphisme parfait (i.e. lisse et propre), alors $f' : X' = X \times_Y Y' \longrightarrow Y'$ est une $W$-équivalence.}

\msp{page 44}
Notons la variante bien naturelle

\textbf{L6 bis} \textit{Comme L6, mais en supposant seulement $Y' \longrightarrow Y$ une} Hot-\textit{fibration.}

\printedpage{28}
% head: A. Grothendieck : Dérivateurs, Ch. VI, Édition des Universités de Montpellier II et Paris VII

Il y a des chances que pratiquement, quand L6 (ou seulement L5) est verifié, L6 bis l'est aussi. % [sic] verifié -> vérifié, missing accent

On va introduire un axiome

\textbf{L7} \ a) \textit{Soit $(f_i : X_i \longrightarrow Y_i)_{i \in I}$ une famille de flèches dans} Cat\textit{, soit $f : \coprod X_i \longrightarrow \coprod Y_i$ la flèche correspondante sur les catégories somme. Pour que $f \in W$, il faut et il suffit que $f_i \in W$ pour tout $i \in I$.}

\hspace*{2em}b) \textit{Soit $f : X \longrightarrow Y$ dans $W$. Si $Y$ est} $0$\textit{-connexe, $X$ aussi.}

Cet axiome équivaut à la conjonction de La ($f \in W \Longrightarrow \pi_0 f$ bijectif) et de la validité du ‘Lemme’ page 23, que je remets ici pour mémoire :

\textbf{Lemme.} \textit{L'axiome L7 précédent sur le localiseur fondamental $W$ équivaut à la condition suivante : $f : X \longrightarrow Y$ est dans $W$ si et seulement si $\pi_0(f)$ est bijectif, et $f$ induit des $W$-équivalences $f_i : X_i \longrightarrow Y_i$ sur les composantes connexes.}

Il me faut revenir encore sur le 1) de la proposition, il n'est pas vrai que ‘c'est’ la condition L4 sur $W$ (que j'admets, ai-je dit), c'est L4 dans le cas particulier où $Y = S$. Réflexion faite, je le vois comme une variante renforcée de L4, qui ne semble pas devoir découler de L4 :

\textbf{L4 bis} \textit{Soit} \begin{tikzcd}[baseline=(s.base), row sep=small, column sep=small]
X \arrow[rr, "f"] \arrow[dr] & & Y \arrow[dl] \\
& |[alias=s]| S &
\end{tikzcd} \textit{dans} Cat\textit{, supposons que $f$ soit une $W$-équivalence locale sur $S$ (i.e. $f/s : X/s \longrightarrow Y/s$ une $W$-équivalence $\forall\ s \in S$), alors $f$ est une $W$-équivalence} (\footnote[29]{équivaut à l'axiome dual, en termes des $s\backslash f$}).

\msp{page 45}
Par la suite, par un ‘localiseur fondamental’, nous entendrons un ensemble
\[
W \subseteq \mathrm{Fl}\,\mathrm{Cat}
\]
satisfaisant (tout au moins) les axioms % [sic] axioms -> axiomes, French spelling
\[
\text{L1, L3 quater, L4 bis, L7,}
\]
i.e.

\begin{itemize}
\item[a)] \textit{Saturation} de $W$ (conditions a) b) c) p. 1)

\item[b)] $W_\infty \subseteq W$ $[\subseteq W_0]$ (\footnote[30]{Je ne suis plus sûr s'il est vraiment utile d'imposer $W_\infty \subseteq W$. Prendre au lieu de ça l'axiome de l'objet final, soit b$'$). Si on veut a) b$'$) c) + L5, alors on aura $W \subseteq W_\infty$. \textbf{NB} a) b) c) + L5 implique $W = W_\infty$, cf. annotation page 43.}) (\footnote[31]{cf. d)}).

\item[c)] Si $f : X \longrightarrow Y$ sur $S$ est localement sur $S$ une $W$-équivalence, alors $f \in W$ (L4 bis)
\end{itemize}

\end{document}
